\section{Пререквизиты}\label{prerequisite}

Для дальнейшей комфортной и плодотворной работы хотелось бы напомнить некоторые основные темы из курса теории вероятностей, которые непременно и многократно пригодятся нам впоследствии. Очевидно, в подобного рода разделе вряд ли удастся пересказать весь необходимый материал, и далее мы будем пользоваться теорией, сильно выходящей за его рамки. Здесь мы не будем даже пояснять, например, что такое математическое ожидание, измеримое отображение и прочее --- не в этом цель. Главное для нас --- подсветить некоторые ключевые моменты сейчас, чтобы не разбирать их на ходу, когда нам приспичит их пустить в ход. Если выражаться метафорически, мы заблаговременно зарядим несколько чеховских ружей, которые будут тут и там выстреливать в течение рассказа. Это отчасти и минус, потому что в данный момент совершенно неясно, зачем мы обсуждаем тот или иной объект из теории вероятностей --- многое покажется нелогичным и непонятно зачем взявшимся. Однако смею Вас заверить, что чтением данных вводных слов зря время потрачено не будет. Итак, приступим.

Отправной точкой в теории вероятности является \textit{вероятностное пространство}: тройка $(\Omega, \mathcal F, \pth[])$, где $\mathcal F$ --- $\sigma$-алгебра на $\Omega$, а $\pth[]$ --- вероятностная мера, заданная на $\mathcal F$. Само по себе оно не представляет особого интереса с практической точки зрения --- это всего лишь математическая модель абстрактного <<чёрного ящика>>, источника случайности. Куда интереснее для нас будет то, что из этого чёрного ящика до нас доносится, а именно случайные величины и векторы. 

Напомним, что \textit{случайным вектором} называется измеримое отображение $\boldxi$ из $(\Omega, \mathcal F)$ в $(\R^k, \mathcal B (\R^k))$, для $k=1$ его обычно называют \textit{случайной величиной}. Измеримость позволяет перенести меру $\pth[]$ на более осязаемое пространство $(\R^k, \mathcal B (\R^k))$ по правилу
$$
\mathsf{P}_{\boldxi}(B) := \pth[](\boldxi \in B), \;\;\; B \in \mathcal B (\R^k);$$
вероятностную меру $\mathsf{P}_{\boldxi}$ называют \textit{распределением случайного вектора $\boldxi$}.

Ввиду своей важности распределения на $\R^k$ хотелось бы как-то легко задавать и описывать. Самым популярным способом является \textit{функция распределения}, которая для распределения $\pth[]$ определяется как
\[
F\colon (x_1, \ldots, x_k) \mapsto \pth[]((-\infty; x_1] \times \ldots \times (-\infty; x_k]).
\]

Широко известно, что по значениям данной функции исходная вероятностная мера определяется однозначно. Более того, если функция $F\colon \R \to  [0; 1]$ обладает не слишком обременительными свойствами (грубо говоря, монотонность, непрерывность справа и какие надо пределы на бесконечности), то существует распределение с функцией распределения~$F$.

\begin{example}
	Для дискретных случайных величин функция распределения представляет собой кусочно-постоянную функцию, претерпевающую разрывы в точках из носителя распределения. Например, для распределения Бернулли $\bernd(p)$, которое определяется как $\pth[](\{1\}) = p$, $\pth[](\{0\}) = 1 - p$, функция распределения имеет вид
	\[
	F(x) = \left\{\begin{aligned}
		&0, \;\; &x < 0;\\
		&1-p,\;\; &x \in [0, 1);\\
		&1, \;\; &x \ge 1.
	\end{aligned}\right.
	\]
	
	Их противоположностью являются распределения с непрерывной функцией распределения, у которых вероятность всякой одной точки равна нулю. Самый простым примером такового является равномерное распределения $\ud[a, b]$, название которого себя полностью оправдывает --- вероятность всякого подотрезка отрезка $[a, b]$ пропорциональна его длине:
	
	\[
	F(x) = \left\{\begin{aligned}
		&0, \;\; &x < a;\\
		&\frac{x-a}{b-a},\;\; &x \in [a, b);\\
		&1, \;\; &x \ge b.
	\end{aligned}\right.
	\]
\end{example}

Отметим полезное свойство функции распределения в одномерном случае, которое позволяет от любого непрерывного распределения перейти к одному и тому же --- равномерному.

\begin{proposition}\label{f_to_uniform}
    Пусть $\xi$ --- случайная величина с непрерывной функцией распределения $F(x)$. Тогда $F(\xi) \sim \ud(0, 1)$.
\end{proposition}

\begin{proof}
    Для $t \in (0, 1)$ определим $F^{-1}(t)$ как $\sup\{x\colon F(x) = t\}$ (множество, по которому берётся супремум, не пусто в силу непрерывности). Тогда из неравенства $F(\xi) \le t$ следует $\xi \le F^{-1}(t)$. В обратную сторону импликация выполняется почти наверное: вероятность события
    \[
    \inf\{x\colon F(x) = t\} < \xi \le \sup\{x\colon F(x) = t\}
    \]
    равна нулю, так как функция распределения $\xi$ одинакова при левой и правой части сего неравенства. Значит,
    \[
    \pth[](F(\xi) \le t) = \pth[](\xi \le F^{-1}(t)) = F(F^{-1}(t)) = t,
    \]
    что есть функция распределения $\ud(0, 1)$. 
\end{proof}

В многомерном случае такое уже не сработает (почему?), да и вообще функции распределения в многомерье становятся не совсем тривиальным объектом для изучения и работы. Но благо есть другой, более простой способ задания распределения --- через плотность. Если нам дана неотрицательная борелевская функция $\rho\colon \R^k \to \R_+$ с единичным интегралом по $\R^k$, то на её основе можно задать распределение следующим образом:
\[
\pth[](B) := \int_{B} \rho(\mathbf x) \, d\mathbf x, \;\;\; B \in \mathcal B(\R^k).
\]

Такие распределения ещё называют \textit{абсолютно непрерывными}. Помимо удобства задания меры мы в подарок получаем удобство подсчёта её характеристик, а именно матожиданий различных (борелевских) функций по ней:
\begin{equation}\label{density_recalc}
	\int_{\R^k} g(\mathbf x) \, \mathsf{P}(d\mathbf x) = \int_{\R^k} g(\mathbf y) \, \rho(\mathbf y ) \, d\mathbf y.
\end{equation}

Минус такого подхода заключается в его узости --- например, дискретные распределения не задашь через интеграл по классической мере Лебега, потому что какую бы плотность мы ни взяли, её интеграл по носителю дискретного распределения будет нулевым. Что же делать? Очень уж хотелось бы и в таком случае определять распределения как интеграл по какой-то одной конкретной мере. Только по какой?

\begin{definition}
    \textit{Считающей мерой} на $\Z^k$ называется функция $\mu\colon \mathcal{B}(\R^k) \to \N_0 \cup \{+\infty\}$, определённая как
    \[
    \mu(B) = \sum_{\mathbf{x} \in \Z^k} I(\mathbf{x} \in B).
    \]
\end{definition}

Ясно, что это $\sigma$-конечная мера, как и классическая мера Лебега на $\R^k$, и она равна числу целочисленных точек, которое попадает в данное множество, откуда собственно и пошло название. Несложно также понять, как считается интеграл произвольной измеримой функции $f$ по этой мере:
\[
\int_{\R^k} f(\mathbf{x})\,\mu(d\mathbf{x}) = \sum_{\mathbf{x} \in \Z^k} f(\mathbf{x}),
\]
если, конечно, этот ряд сходится абсолютно. Такое представление позволяет записывать матожидание от дискретных случайных векторов $\xi\colon \Omega \to \Z^k$ через интеграл с плотностью по считающей мере, обобщая равенство \eqref{density_recalc}:
\begin{equation}\label{density_recalc_discr}
	\me[] g(\xi) = \sum_{\mathbf{x} \in \Z^k} g(\mathbf{x}) \cdot \pth[](\xi = \mathbf{x}) = \int_{\R^k} g(\mathbf{x}) \cdot \rho(\mathbf{x})\,\mu(d\mathbf{x}),
\end{equation}
где $\rho(\mathbf{x}) = \pth[](\xi = \mathbf{x})$ --- \textit{плотность по мере $\mu$}. Отныне под словами <<плотность по мере $\mu$>> мы будем подразумевать либо обычную плотность, которая у нас была ранее, по классической мере Лебега, либо дискретную по считающей мере. Семейства распределений, у которых есть плотность по одной и той же мере, мы будем называть \textit{доминируемым}.

\begin{example}
	В примере выше семейство равномерных распределений на отрезке $\{\ud(a, b)\colon a < b\}$ доминируется классической мерой Лебега на прямой с плотностью $\rho_{\ud(a, b)}(x) = \frac{1}{b-a} \cdot I\{x \in (a, b)\}$ --- её постоянство на носителе $(a, b)$ служит ещё одним аргументом в пользу такого названия. 
		
	Семейство бернуллиевских распределений $\{\bernd(p)\colon p \in (0, 1)\}$ доминируемо по считающей мере с плотностью
	\[
	\rho_{\bernd(p)}(x) = \left\{
		\begin{aligned}
			&1-p, & x = 0,\\
			&p,   & x = 1,\\
			&0,   & x \not\in \{0, 1\}.
		\end{aligned}
	\right.
	\]
	
	В дальнейшем нам будет удобно записывать эту плотность в более коротком виде, $\rho_{\bernd(p)}(x) = p^x (1-p)^{1-x} \cdot I(x \in \{0, 1\})$.
\end{example}

{
\footnotesize Данные соображения можно распространить на случай произвольной $\sigma$-конечной меры $\mu$ с помощью теоремы Радона-Никодима.

\begin{definition}
     Пусть $\mu$ --- некоторая $\sigma$-конечная мера на $\R^k$. Семейство вероятностных мер $\mathcal{P}$ называется \textit{доминируемым} относительно меры $\mu$, если $\forall \mathsf{P} \in \mathcal{P}\colon \mathsf{P} \ll \mu$ (напомним, что $\nu \ll \mu$, если $\forall B \in \mathcal{B}(\R^k)\colon \mu(B) = 0 \Rightarrow \nu(B) = 0$). Производную Радона-Никодима $\frac{d\mathsf{P}}{d\mu}$ называют \textit{обобщённой плотностью}.
\end{definition}

Несложно показать, что и в этом случае выполнена \textit{формула пересчёта}, аналогичная \eqref{density_recalc}. Она позволяет перейти от интеграла по неизвестной мере к интегралу плотности по известной мере, которая и доминирует семейство:
\[
\int_{\R^k} f(\mathbf{x})\,\mathsf{P}(d\mathbf{x}) = \int_{\R^k} f(\mathbf{x}) \cdot \frac{d\mathsf{P}}{d\mu}\,\mu(d\mathbf{x}).
\]
}

На этом моменте уже можно изучать различные распределения, как они меняются у случайных векторов при действии на них отображений и т.д., однако довольно скоро вы поймёте, что всё хорошее быстро закончится. Надежды получить распределение хоть сколь-нибудь сложных объектов в явном виде рушатся при взятии первого же интеграла, чего только формула свёртки стоит. При этом на практике мы будем встречаться с массой подобных примеров, отчего возникает необходимость в разработке обходных путей. Вдруг можно считать распределения не явно, а хотя бы приближённо?

Для начала вспомним, а что вообще значит <<случайный вектор $\boldxi_n$ близок к вектору $\boldxi$>>. Есть как минимум три способа понимать эту фразу:

\begin{itemize}
    \item сходимость почти наверное ($\boldxi_n \stackrel{\text{п.н.}}{\to} \boldxi$), когда $\mathsf{P}(\{\omega\colon \boldxi_n(\omega) \stackrel{\text{п.н.}}{\to} \boldxi(\omega)\}) = 1$;
    \item сходимость по вероятности ($\boldxi_n \stackrel{\pth[]}{\to} \boldxi$), когда для любого $\epsilon > 0$ выполнено $\mathsf{P}(\|\boldxi_n - \boldxi\| > \epsilon) \to 0$;
    \item слабая сходимость или сходимость по распределению ($\boldxi_n \stackrel{d}{\to} \boldxi$), когда $F_{\boldxi_n}(\mathbf{x}) \to F_{\boldxi}(\mathbf{x})$ для всякой точки $\mathbf{x}$, в которой $F_{\boldxi}$ непрерывна.
\end{itemize}

Первый вид сходимости является естественным обобщением поточечной сходимости, при котором нам неважно, что происходит на событии вероятности нуль. Второй исторически возник из закона больших чисел. Третий же полезен тем, что он понимает близость случайных векторов через близость их распределений. Действительно, в определении используются только $F_{\boldxi}$ и $F_{\boldxi_n}$, поэтому эту сходимость можно понимать ещё и как сходимость вероятностных мер. В частности, мы будем часто писать, например, $\boldxi_n  \stackrel{d}{\to} \mathcal{N}(0, 1)$, подразумевая слабую сходимость распределений векторов слева к распределению справа.

Широко известно, что из сходимости почти наверное следует сходимость по вероятности, а из последней -- сходимость по распределению. Хотелось бы поиметь больший арсенал для получения одних сходимостей из других. Важнейшим таким источником является

\begin{theorem}[label=clt]{о наследовании сходимости}{}
    Пусть $g\colon \R^m \to \R^k$ --- борелевская функция, непрерывная на множестве $B$ таком, что $\pth[](\boldxi \in B) = 1$. Тогда из сходимости $\boldxi_n$ к $\boldxi$ (в любом из смыслов: п.н., по вероятности или по распределению) следует сходимость $g(\boldxi_n)$ к $g(\boldxi)$ в том же смысле.
\end{theorem}

Другой важный аспект --- связь сходимостей случайных векторов $\boldxi_n \to \boldxi$ и его компонент $\boldxi_n^{(i)} \to \boldxi^{(i)}$. Легко показать, что в случаях сходимости п.н. и по вероятности эти вещи равносильны, но для слабой сходимости импликация верна только вправо, ведь слабая сходимость компонент вектора ничего не говорит об их совместном распределении. К счастью, если два простых достаточных условия, когда из слабой сходимости $\boldxi_n^{(i)} \stackrel{d}{\to} \boldxi^{(i)}$ для всех $i$ следует $\boldxi_n \stackrel{d}{\to} \boldxi$:

	\begin{enumerate}
		\item Если компоненты векторов независимы (см. задачу \ref{weak_conv_while_ind});
		\item Если один из пределов является константой.
	\end{enumerate}

Последнее утверждение часто выделяют отдельно как

\begin{theorem}{Слуцкий}{}
    Пусть $\xi_n \stackrel{d}{\to} \xi$ и $\eta_n \stackrel{d}{\to} c$, где $c$ --- некоторая константа. Тогда $(\xi_n, \eta_n) \stackrel{d}{\to} (\xi, c)$.
\end{theorem}

В частности, применяя теорему \ref{clt}, мы делаем вывод, что имеют место сходимости $\xi_n + \eta_n \stackrel{d}{\longrightarrow} \xi + c$ и $\xi_n \cdot \eta_n \stackrel{d}{\longrightarrow} \xi \cdot c$, то есть слабые сходимости можно складывать и умножать, если один из пределов постоянен. Этот факт мы для краткости будем называть\textit{ леммой Слуцкого}.

Напоследок вспомним, откуда мы можем черпать различные нетривиальные сходимости. В курсе теории вероятностей выделяют два важнейших предельных результата.

Первый из них, закон больших чисел (ЗБЧ) утверждает, что среднее арифметическое независимых одинаково распределённых величин с ростом их числа приближается к теоретическому среднему, то есть матожиданию. Первоначально теорема говорила о сходимости по вероятности, однако со временем человечество поняло, что при условиях ниже имеет место даже сходимость почти наверное --- этот результат принято называть усиленным законом больших чисел (УЗБЧ).

\begin{theorem}{(усиленный) закон больших чисел}{}
    Пусть $X_1, X_2, \ldots$ --- независимые одинаково распределённые случайные величины с конечным матожиданием $\me[] X_1 = a$. Тогда
    \[
    \frac{X_1 + \ldots + X_n}{n} \stackrel{\pth[]}{\longrightarrow} a \;\; \qquad \left(\frac{X_1 + \ldots + X_n}{n} \stackrel{\text{п.н.}}{\longrightarrow} a \right).
    \]
\end{theorem}

Имеется множество обобщений данного утверждения, отказывающихся от условий одинаковой распределённости и независимости взамен на существование б\'{о}льших моментов, но обычно нам будет хватать теоремы выше.

Закон больших чисел говорит лишь о самом факте сходимости, но умалчивает о её скорости и о том, как устроено отклонение среднего арифметического от $a$. Ответ на это даёт второй предельный результат и, пожалуй, самая знаменитая теорема теории вероятностей:

\begin{theorem}{центральная предельная теорема}{}
    Пусть $X_1, X_2, \ldots$ --- независимые одинаково распределённые случайные величины с матожиданием $\me[] X_1 = a$ и конечной ненулевой дисперсией $\va[] X_1 = \sigma^2$. Тогда
    \[
    \sqrt{n}\left(\frac{\sum_{i=1}^n X_i}{n} - a\right) \stackrel{d}{\longrightarrow} \mathcal{N}(0, \sigma^2), \quad \text{или, что то же самое,}\quad \frac{\sum_{i=1}^n X_i - na}{\sigma\sqrt{n}} \stackrel{d}{\longrightarrow} \mathcal{N}(0, 1).
    \]
\end{theorem}

Первое, что говорит нам теорема --- сходимость в УЗБЧ имеет порядок $1 / \sqrt{n}$. Может быть не совсем понятно, как говорить о <<скорости сходимости>> для последовательности случайной, но достаточно использовать такие соображения:
\begin{itemize}
	\item если бы мы домножили разность $\frac{\sum_{i=1}^n X_i}{n} - a$ на что-то менее быстро растущее, чем $\sqrt{n}$, то по лемме Слуцкого предел будет нулевым;
	\item если бы мы домножили разность на что-то более быстро растущее, чем $\sqrt{n}$, то последовательность не просто не будет иметь предела по распределению, но и в каком-то смысле <<стремиться>> к бесконечности.
\end{itemize}

Второе --- теорема даёт нам ту самую аппроксимацию, которую можно использовать вместо настоящего распределения среднего. Неформально эту сходимость можно понимать так: раз уж величина $\sqrt{n}\left(\frac{\sum_{i=1}^n X_i}{n} - a\right)$ <<почти что>> нормально распределена с нулевым средним и дисперсией $\sigma^2$, то перенося множитель $\sqrt{n}$ и слагаемое $a$ вправо, можно считать, что
\[
\frac{\sum_{i=1}^n X_i}{n} \stackrel{d}{\approx} \mathcal{N}\left(a, \frac{\sigma^2}{n}\right).
\]

Обычно именно это имеют в виду, когда пользуются ЦПТ, однако сходимость выше позволяет эту, вообще говоря, безграмотную запись формализовать.


Обе теоремы имеют многомерные аналоги. Для (У)ЗБЧ её даже формулировать стыдно, ибо всё очевидно следует из связи сходимости векторов и их компонент. Для ЦПТ же приведём явную формулировку, ибо в ней появляется важный объект.

\begin{theorem}{многомерная центральная предельная теорема}{}
    Пусть $\boldxi_1, \boldxi_2, \ldots$ --- независимые одинаково распределённые случайные $k$-мерные векторы с вектором средних $\me[] \boldxi_1 = \mathbf a$ и конечной ковариационной матрицей $\Sigma = \va[] \boldxi_1$. Тогда
    \[
    \sqrt{n}\left(\frac{1}{n}\sum_{i=1}^n \boldxi_i - \mathbf a\right) \stackrel{d}{\longrightarrow} \mathcal{N}(\mathbf 0, \Sigma).
    \]
\end{theorem}

В данном случае предел является обобщением нормального распределения, а именно гауссовским вектором. Без лишних подробностей, одно из его определений --- это аффинное преобразование вектора из независимых случайных величин с распределением $\mathcal{N}(0, 1)$, у которого вектор средних и ковариационная матрица равны $\mathbf a$ и $\Sigma$ соответственно. Как и в одномерном случае, распределение полностью описывается этими двумя параметрами. Из важных свойств гауссовского вектора выделим следующее известное

\begin{proposition}[критерий независимости компонент]\label{cov_and_ind_are_eq}
    Пусть дан $n$-мерный гауссовский вектор
$$
\boldxi = (\xi_{1}^1, \ldots, \xi_{k_1}^1, \xi_{1}^2, \ldots, \xi_{k_2}^2, \ldots, \xi_{1}^s, \ldots, \xi_{k_s}^s), \;\;\; n = k_1 + \ldots + k_s.
$$
Тогда его компоненты $\boldxi_i = (\xi_{1}^i, \ldots, \xi_{k_i}^i)$, $i=1, \ldots, s$, независимы в совокупности тогда и только когда, когда матрица ковариаций блочно-диагональная, а именно
$$
\forall i, j = 1, \ldots, s; i \ne j,\;\forall m = 1, \ldots, k_i;  \forall l = 1, \ldots, k_j\colon \cov(\xi^i_{m}, \xi^j_{l}) = 0.
$$
\end{proposition}

В частном случае, когда $k_1 = \ldots = k_s = 1$ получаем простое утверждение: компоненты гауссовского вектора независимы $\Longleftrightarrow$ они не коррелированы.

\subsection*{Задачи}

\begin{problem}\label{weak_conv_while_ind}
	Пусть $\boldxi$ и $(\boldxi_n\colon n \in \N)$ --- случайные $k$-мерные векторы, у каждого из которых компоненты независимы в совокупности, причём $\boldxi_n^{(i)} \stackrel{d}{\to} \boldxi^{(i)}$ для $i \in \{1, \ldots, k\}$. Докажите, что $\boldxi_n \stackrel{d}{\to} \boldxi$.
\end{problem}

\begin{problem}
	Выведете многомерную ЦПТ из её одномерного случая.
\end{problem}

